\documentclass[../main.tex]{subfiles}
\graphicspath{{\subfix{../images/}}}
\begin{document}

\begin{figure}[htb!]
\centering
  \includegraphics[width=7cm]{fat-foods}
  \caption{Foods rich in fats~\cite{FatFoods}}
\end{figure}

Fats are the most important fuels for the human body, next to carbohydrates.\index{fat}
They have about 9 kcal (\SI{37}{\kilo\joule} per gram (\SI{255}{\kilo\joule}/ \SI{1050}{kJ} per oz.), a bit more than double the energy content than carbohydrates and fats.
Fats also have a big signification as building blocks for cells and organs.
The brain for instance consists of 50\% of fat, the white brain matter even 70\%.
But fats are also very important when it's about protecting the organism,
either as insulation against heat and cold or as a cushion for organs and the nervous system.
One of the main constituents of fats are the fatty acids\index{fatty acid}.
One part of these fatty acids, like for instance omega--3 and omega--6 fatty acids\index{fatty acid!omega--3}\index{fatty acid!omega--6}
can't be build by the human body, therefor they are essential to us.
Fatty acids are contained in both, in plant based and animal based foods.

Fats or lipids have the following characteristics:
\begin{itemize}
\item important fuel for our cells, as body fat also as an energy reserve
\item as body fat, it protects the organism of heat and cold, and the organs and nervous system of mechanical shocks
\item they are carrier of the fat--soluble vitamins A, D, E and K\index{vitamin!fat--soluble}
\item in foods, fats are popular as carriers of taste. On top of that
  they cause a good feeling of saturation and prolong the passage of the foods through the colon
\item fats are an important ingredient for the cell membranes\index{cell membranes}
\item certain fats are contributing to the mechanisms of the immune system\index{immune system}
\item they are the building blocks for hormones and vitamins\index{hormones} 
\end{itemize}

\subsection{Structure of Fats}

In nutrition, fats appear in solid form as fats or in liquid form as oils.
Nutritional fats (triglycerides)\index{triglycerides} are compounds consisting of glycerol and three fatty acids, each.
For the fatty acids, we distinguish  between saturated,
monounsaturated and polyunsaturated fatty acids.\index{fatty acid!saturated}\index{fatty acid!unsaturated}

\vspace{3mm}
{\centering
  \includegraphics[width=11cm]{FattyAcids}
}

\begin{description}
\item[Saturated Fatty Acids] can be build up by human body.
  Certain saturated fatty acids increase the concentration of cholesterol\index{cholesterol} in the blood.
\item[Some Polyunsaturated Fatty Acids] can't be build up by our body and are therefore essential\index{fatty acid!essential}.
  Among those are omega--6 fatty acids\index{fatty acid!omega--6} linoleic acid\index{fatty acid!linoleic} and the
  omega--3 fatty acids\index{fatty acid!omega--3} alpha--linolenic acid\index{fatty acid!alpha--linolenic} and its derivatives
  eicosapentaenoic acid (EPA)\index{EPA} and docosahexaenoic acid (DHA).\index{DHA}
  Contrary to saturated fatty acids, polyunsaturated fatty acids can lower the cholesterol levels in the blood.
\end{description}

\begin{table}[htb!]
  \centering
  \begin{tabular}{l|ll}
    \textbf{Fatty Acid} & \multicolumn{2}{c}{\textbf{Foods}} \\
                        & \textbf{animal base} & \textbf{plant based} \\
    \hline
    Saturated & meat, sausage, cheese, cream & coco and palm oils, \\
                        & lard, shortening, suet \\
    Monounsaturated & goose fat, eggs & olive, canola, and soy oils, \\
                        & & nuts, avocado \\
    Polyunsaturated & \multicolumn{2}{c}{\textit{Omega--3:}} \\
                        & fatty ocean fish (herring,  & linen, soy , hemp and \\
                        & salmon, mackerel), anchovy & walnut oils, walnuts, algae,\\
                        & sardine, venison & dark green leafy greens \\
  & \multicolumn{2}{c}{\textit{Omega--6:}} \\  
                        & meat, butter, milk and & sunflower, thistle, \\
                        & milk products & pumpkin seed, corn seed and \\
                        & & grape seed oils \\

  \end{tabular}
  \caption{Foods and the fatty acids they contain}
\end{table}

\begin{table}[htb!]
  \centering
  \begin{tabular}{lrrrr}
    \textbf{Fat intake (daily)} \\
    \multicolumn{5}{l}{Guide values in \% of the daily calorie intake*}\\
    \hline
    Children & Girls & 30--35\% & Boys & 30--35\% \\
    Adolescents & Girls & 30\% & Boys & 30\% \\
    Adults & Women & 30\% & Men & 30\% \\
    Pregnant women & after 4th month & 30--35\%\\
    Breastfeeding women & & 30--35\%\\
    \hline
    \hline
    \multicolumn{5}{l}{Depending of caloric need, a fat ratio of 30\% corresponds to:}\\
    \hline
    2000 kcal & 65 g fat & 2\sfrac{1}{4} oz\\
    3000 kcal & 100 g fat & 3\sfrac{1}{2} oz\\
    \hline
         \footnotesize{*Reference values DACH 2000}
  \end{tabular}
  \caption{How much fat does a human being need?}
\end{table}

Beside the total fat intake is the ratio of the fatty acids of importance.

\begin{table}[htb!]
  \centering
  \begin{tabular}{ll}
    \textbf{Fatty acid} & \textbf{Daily intake} \\
    \hline
    Saturated & \sfrac{1}{3} of the fat intake or max. 10\% of the caloric intake \\
    Monounsaturated & \sfrac{1}{3} of the fat intake or min. 10\% of the caloric intake \\
    Polyunsaturated & \sfrac{1}{3} of the fat intake or 7 -- 10\% of the caloric intake \\
                       & of that Omega--6 to Omega--3 in a ratio of 5:1\\
    Trans fatty acids & less than 1\% of the caloric intake \\
    Cholesterol & 300 mg \\
  \end{tabular}
  \caption{An optimal ratio of fatty acids}
\end{table}

\subsection{Omega--3, the Special Fatty Acid}
The omega--3 fatty acids are an uttermost important constituent of a healthy way of live.
Especially the omega--3 fatty acid docosahexaenoic acid (DHA) is very important for a healthy nutrition.

Newer scientific research showed, that omega--3 fatty acids have following positive effects on the human organism:
\begin{itemize}
\item activate genes for the combustion of fats\index{fat combustion}
\item lower the blood pressure\index{effect!blood pressure!lower}\index{symptom!blood pressure high}
\item lower the blood fat content (cholesterol)\index{effect!cholesterol!lower}\index{symptom!cholesterol}
\item lower the risk for thrombosis (blood clots) by inhibiting the coagulation
  of the blood and by improving the flowing capacity of the blood\index{effect!thrombosis risk!lower}\index{symptom!thrombosis}
  \item help with allergies and auto--immune diseases\index{smyptom!allergies}\index{symptom!auto--immune diseases}
  \item are important for the development of the brain
  \item have a positive influence with ADHS, depression, stress, aggression,
    cancer, Alzheimer\index{symptom!Alzheimer} and many other health concerns.\index{symptom!ADHS}\index{symptom!depression}\index{symptom!stress}\index{symptom!aggression}\index{symptom!cancer}
\end{itemize}

Omega--3 fatty acids aren't always sufficiently contained in our food.
Many developed countries, like middle of Europe of America, have very little fish in their food and therefore a general lack of intake of DHA.
The valuable DHA is mainly coming form fish.
It is therefore recommended to eat at least twice a week fish or to supplement\index{supplement!omega--3} the nutrition with omega--3 in the form fish oil capsules.
But fish is not to everybody's taste and furthermore a valuable resource.
We often read that the fish population in the oceans are dwindling.

The Swiss company Allsan succeeded in producing a 100\% plant based replacement to the fish oil products.
The product Allsan Omega3Vega is gained from a specially breed algae and contains foremost the fatty acid DHA.
This DHA is identical to DHA gained from fish in respect to quality.

 
\vspace{5mm}
\noindent
\begin{fminipage}{\textwidth}
  \textbf{Profile Fats}
  \begin{itemize}
  \item fats = lipids
  \item the building blocks of fats are fatty acids
  \item 1 g fats delivers 9kcal (37kJ) energy, 1 oz delivers 255 kJ or 1050 kJ
  \item recommended daily intake of fats, according to DACH (19--65 year old): 30\% of the daily caloric intake
  \item distinguish saturated, monounsaturated and polyunsaturated fatty acids
  \item omega--6 and omega--3 are essential and should be supplied to the body through foods in the ratio 5:1
  \item plant based sources of fats (to be preferred): plant oils, nuts, avocado, linen seeds, algae, dark green leafy greens
  \item animal based sources of fat: fish, venison, eggs, milk and milk products, cheese
    \item omega--3 is often recommended to supplement food --- Allsan Omega3Vega (purely plant based) as an alternative to the fish oil capsules 
  \end{itemize}
\end{fminipage}

\end{document}
